\begin{figure}[!tb]
    \centering
    \includegraphics[width=\linewidth]{figures/overview-rebuttal.pdf}
    \caption{
    This survey covers the diverse ways physics knowledge can be incorporated into ML across datasets $\gD$, objectives $\gL$ and models $\vu_\vtheta$.
    Physics-informed models are architectures and prediction techniques informed by domain knowledge, covered in sections \ref{sec:node}, \ref{sec:equations}, and \ref{sec:lvm}.
    Learning arbitrary models with physics-informed losses is discussed in sections \ref{sec:pinns} and \ref{sec:sysid}.
    Data augmentation can also incorporate physics knowledge to encode invariances (section \ref{sec:aug}).
    Data-driven learning across multiple datasets and experiments includes neural operators (\secref{sec:no}), multi-task learning (section \ref{sec:multitask}) and meta-learning. %(Section \ref{sec:meta}), and has been used for both physics-informed and -uninformed models and losses.
    Finally, contextual data is another data-driven approach to incorporate additional domain knowledge, as seen in operator learning (\secref{sec:no}) and neural processes (\secref{sec:np}).
    }
    \label{fig:overview}
\end{figure}
